\section{Experiments}

\subsection{Setup}

\noindent\textbf{Dataset.} COD10K-v3~\cite{fan2020cod10k}: 3,040 camouflaged training images, 2,026 camouflaged test images across 78 categories (aquatic, terrestrial, flying, amphibian).

\noindent\textbf{Implementation.} SAM ViT-H backbone. Training images augmented with horizontal flips (6,080 total). Decoder training: $<$30 min on A100 with pre-computed embeddings. LLPM+decoder training: $\sim$12 hrs (live encoder pass). Decoder checkpoint: 16MB.

\subsection{Main Results}

Table~\ref{tab:main} presents results across all prompt strategies.

\begin{table*}[t]
\centering
\caption{\textbf{Main results} on COD10K camouflaged test set (2,026 images). Best per-column in \textbf{bold}. Decoder-only fine-tuning improves all metrics across all prompt types; LLPM adds further gains on edge prompts.}
\label{tab:main}
\small
\begin{tabular}{llccccccc}
\toprule
Model & Prompt Strategy & mIoU & Dice & $S_\alpha$ & $E_\phi$ & $F_\beta^w$ & MAE$\downarrow$ & Bnd.~F1 \\
\midrule
Base SAM & Center-of-Mass (1pt) & 0.572 & 0.654 & 0.743 & 0.830 & 0.706 & 0.068 & 0.502 \\
Base SAM & Edge (1pt) & 0.205 & 0.251 & 0.481 & 0.597 & 0.241 & 0.163 & 0.281 \\
Base SAM & Multi-Pt Grid (4pt) & 0.682 & 0.761 & 0.812 & 0.897 & 0.802 & 0.046 & 0.600 \\
Base SAM & Multi-Pt Random (3pt) & 0.726 & 0.812 & 0.830 & 0.928 & 0.831 & 0.045 & 0.648 \\
\midrule
Decoder-only & Center-of-Mass (1pt) & 0.738 & 0.819 & 0.847 & 0.942 & 0.819 & 0.028 & 0.702 \\
Decoder-only & Edge (1pt) & 0.694 & 0.779 & 0.810 & 0.922 & 0.750 & 0.034 & 0.666 \\
Decoder-only & Multi-Pt Grid (4pt) & \textbf{0.756} & \textbf{0.835} & \textbf{0.861} & \textbf{0.951} & \textbf{0.841} & \textbf{0.024} & \textbf{0.720} \\
Decoder-only & Multi-Pt Random (3pt) & \textbf{0.774} & \textbf{0.855} & \textbf{0.874} & \textbf{0.959} & \textbf{0.862} & \textbf{0.022} & \textbf{0.739} \\
\midrule
LLPM + Decoder & Center-of-Mass (1pt) & \textbf{0.740} & \textbf{0.820} & \textbf{0.847} & \textbf{0.942} & \textbf{0.820} & \textbf{0.028} & \textbf{0.702} \\
LLPM + Decoder & Edge (1pt) & \textbf{0.719} & \textbf{0.802} & \textbf{0.827} & \textbf{0.922} & \textbf{0.778} & \textbf{0.029} & \textbf{0.690} \\
LLPM + Decoder & Multi-Pt Grid (4pt) & 0.755 & 0.833 & 0.856 & 0.951 & 0.833 & 0.025 & 0.718 \\
LLPM + Decoder & Multi-Pt Random (3pt) & 0.772 & 0.851 & 0.866 & 0.959 & 0.848 & 0.022 & 0.735 \\
\bottomrule
\end{tabular}
\end{table*}

Decoder-only fine-tuning improves mIoU on all four prompt types by +4.8 to +48.9 absolute points. The LLPM provides an additional +2.5 mIoU on edge prompts, specifically the scenario where boundary signals matter most.

\subsection{Prompt Robustness}

\begin{table}[t]
\centering
\caption{\textbf{Prompt robustness.} Base SAM is highly sensitive to prompt quality (mIoU 0.205--0.726); our decoder narrows this to 0.694--0.774.}
\label{tab:improvement}
\small
\begin{tabular}{lcccc}
\toprule
Prompt Type & Base & Ours & $\Delta$ & Rel.~$\Delta$ \\
\midrule
Center-of-Mass & 0.572 & 0.738 & +0.166 & +28.9\% \\
Edge (Single) & 0.205 & 0.694 & +0.489 & \textbf{+238.1\%} \\
Multi-Pt Grid & 0.682 & 0.756 & +0.074 & +10.9\% \\
Multi-Pt Random & 0.726 & 0.774 & +0.048 & +6.6\% \\
\bottomrule
\end{tabular}
\end{table}

The headline result: base SAM's mIoU varies 3.5$\times$ across prompt types (0.205 to 0.726), while the specialized decoder varies only 1.1$\times$ (0.694 to 0.774). Edge prompts---where the prompt sits at the boundary between camouflaged object and background---show the largest absolute gain (+48.9 mIoU). This demonstrates that decoder specialization teaches the model to resolve boundary ambiguity rather than memorize prompt-mask associations.

\subsection{Per-Category Analysis}

\begin{table}[t]
\centering
\caption{\textbf{Per-category mIoU} (Center-of-Mass prompt). Terrestrial animals---the hardest category---show the largest relative gain.}
\label{tab:category}
\small
\begin{tabular}{lcccr}
\toprule
Category & $N$ & Base & Ours & Rel.~Gain \\
\midrule
Amphibian & 124 & 0.680 & 0.825 & +21.4\% \\
Aquatic & 474 & 0.563 & 0.720 & +27.9\% \\
Flying & 714 & 0.610 & 0.767 & +25.6\% \\
Terrestrial & 699 & 0.519 & 0.703 & \textbf{+35.5\%} \\
\bottomrule
\end{tabular}
\end{table}

Terrestrial animals (insects on bark, lizards on rocks) show the largest gain (+35.5\%), consistent with having the most challenging camouflage. Aquatic animals (+27.9\%) also benefit substantially.

\subsection{Error Analysis}
\label{sec:error_analysis}

To understand \emph{where} our model still fails, we analyze the worst-performing images (bottom quartile by IoU) and categorize them using image metadata computed during evaluation.

\noindent\textbf{Failure categories.} We define five automatic categories:
\begin{itemize}
    \item \textbf{Dark}: image brightness $<$ 80 (out of 255)
    \item \textbf{Tiny}: foreground pixels $<$ 1\% of image area
    \item \textbf{Thin/elongated}: GT bounding box aspect ratio $>$ 5:1
    \item \textbf{Disjoint}: multiple connected components in GT mask
    \item \textbf{General}: none of the above
\end{itemize}

\begin{table}[t]
\centering
\caption{\textbf{Error analysis.} Mean IoU on the 400 worst-performing images, grouped by failure category. LLPM improves on all categories except the rare thin/elongated case ($N\!=\!2$).}
\label{tab:error_analysis}
\small
\begin{tabular}{lrccc}
\toprule
Category & $N$ & Base & Spec. & LLPM \\
\midrule
Disjoint & 175 & .213 & .303 & \textbf{.371} \\
Dark & 104 & .316 & .326 & \textbf{.366} \\
General & 65 & .352 & .342 & \textbf{.413} \\
Tiny & 54 & .289 & .261 & \textbf{.292} \\
Thin/elongated & 2 & .004 & \textbf{.346} & .222 \\
\bottomrule
\end{tabular}
\end{table}

\noindent\textbf{Key findings.}
Failures are structured, not random (Table~\ref{tab:error_analysis}). Disjoint objects (multiple connected components, 44\% of failures) and dark images (26\%) dominate. The LLPM improves over decoder-only on all major failure categories: +6.8 mIoU on disjoint objects, +7.1 on general failures, and +4.0 on dark images. This demonstrates that the LLPM's boundary enhancement benefits extend beyond edge prompts to the model's hardest cases overall. Tiny objects ($<$1\% foreground pixels) remain the most challenging category (mean IoU 0.29), suggesting that scale-aware augmentation or multi-scale features are needed for further improvement.

\subsection{Comparison with Existing Methods}

\begin{table}[t]
\centering
\caption{\textbf{Comparison with COD methods} on COD10K test set. Our method uses GT prompts; dedicated methods are fully automatic.}
\label{tab:comparison}
\small
\begin{tabular}{llcccc}
\toprule
Method & Venue & $S_\alpha$ & $E_\phi$ & $F_\beta^w$ & MAE$\downarrow$ \\
\midrule
SINet~\cite{fan2020cod10k} & CVPR'20 & .776 & .864 & .631 & .043 \\
PFNet~\cite{mei2021pfnet} & CVPR'21 & .800 & .877 & .660 & .040 \\
SINet-V2~\cite{fan2022sinetv2} & TPAMI'22 & .815 & .887 & .680 & .037 \\
SegMaR~\cite{jia2022segmar} & CVPR'22 & .833 & .899 & .724 & .034 \\
ZoomNet~\cite{pang2022zoomnet} & CVPR'22 & .838 & .911 & .729 & .029 \\
\midrule
SAM-Adpt.~\cite{chen2023samadapter} & ICCVW'23 & .883 & .918 & .801 & .025 \\
COMPr.~\cite{hao2024comprompter} & SCIS'24 & \textbf{.889} & .949 & .821 & \textbf{.023} \\
\midrule
\textbf{Ours} (Grid) & --- & .861 & \textbf{.951} & \textbf{.841} & .024 \\
\bottomrule
\end{tabular}
\end{table}

Our decoder-only approach surpasses all dedicated COD architectures on $E_\phi$ and $F_\beta^w$ despite requiring no domain-specific design. Methods that modify the encoder (SAM-Adapter, COMPrompter) achieve higher $S_\alpha$---consistent with our thesis that encoder modification and decoder fine-tuning address complementary bottlenecks. Note: our method uses GT prompts; dedicated methods are fully automatic.
